%% linattn-chunkwise.tex -- the compromise between the two previous forms: the
%% scan runs over chunks rather than tokens, and inside a chunk the masked
%% matrix form runs in parallel.
%%
%% The three operands keep the colours of the other two figures, so a chunk
%% reads as a miniature of the parallel form.
%%
%% Requires texfig.sty. Input inside a figure* environment:
%%   \input{figures/linattn-chunkwise}
\providecommand{\tffigwidth}{\linewidth}
\resizebox{\tffigwidth}{!}{%
\begin{tikzpicture}[font=]

% ===================================================================
% The chunks, each holding its own Q, K and V
% ===================================================================
\foreach \name/\x/\s in {ca/3.9/1, cb/9.9/2, cc/17.1/M} {
  \node[tflabel, font=\small\bfseries] (\name h) at (\x,1.15) {Chunk $\s$};
  \draw[tfglyph=tfRed]   (\x-0.95,0.05) ellipse (0.17 and 0.52);
  \draw[tfglyph=tfBlue]  (\x,0.05)      ellipse (0.45 and 0.32);
  \draw[tfglyph=tfGreen] (\x+0.95,0.05) ellipse (0.17 and 0.52);
  \node[tflabelc=tfRed, font=\scriptsize]   (\name q) at (\x-0.95,-0.85)
        {$\mathbf{Q}_{\s}$};
  \node[tflabelc=tfBlue, font=\scriptsize]  (\name k) at (\x,-0.85)
        {$\mathbf{K}_{\s}^\top$};
  \node[tflabelc=tfGreen, font=\scriptsize] (\name v) at (\x+0.95,-0.85)
        {$\mathbf{V}_{\s}$};
  \begin{scope}[on background layer]
    \node[tfcard, inner sep=0.3cm, fit=(\name h)(\name q)(\name k)(\name v)]
         (\name) {};
  \end{scope}
}

% ===================================================================
% The state handed from one chunk to the next, pinned to the line of the
% chunk cards so every arrow along the scan is horizontal
% ===================================================================
\node[tfdot=tfGray, minimum size=1.2cm] (s0) at (\tfrow{0.9}{ca.west})  {$\mathbf{S}^{(0)}$};
\node[tfdot=tfGray, minimum size=1.2cm] (s1) at (\tfrow{6.9}{ca.west})  {$\mathbf{S}^{(1)}$};
\node[tfdot=tfGray, minimum size=1.2cm] (s2) at (\tfrow{12.9}{ca.west}) {$\mathbf{S}^{(2)}$};
\node[tflabel, font=\large]                (dots) at (\tfrow{15.0}{ca.west}) {$\cdots$};
\node[tfdot=tfGray, minimum size=1.2cm] (sm) at (\tfrow{20.1}{ca.west}) {$\mathbf{S}^{(M)}$};

\draw[tfarrowc=tfGreen] (s0.east) -- (ca.west);
\draw[tfarrowc=tfGreen] (ca.east) -- (s1.west);
\draw[tfarrowc=tfGreen] (s1.east) -- (cb.west);
\draw[tfarrowc=tfGreen] (cb.east) -- (s2.west);
\draw[tfarrowc=tfGreen] (s2.east) -- (dots.west);
\draw[tfarrowc=tfGreen] (dots.east) -- (cc.west);
\draw[tfarrowc=tfGreen] (cc.east) -- (sm.west);

% ===================================================================
% What each chunk emits on its own
% ===================================================================
\foreach \name/\s in {ca/1, cb/2, cc/M} {
  \node[tfghost] (\name o) at (\tfcol{\name.south}{-3.0}) {$\mathbf{O}^{(\s)}$};
  \draw[tfarrowc=tfRed] (\name.south) -- (\name o.north);
}

% ===================================================================
% The two terms the chunking splits the output into
% ===================================================================
\node[tfformula] (eq) at (8.4,-5.9)
      {$\mathbf{O}^{(k)} = \underbrace{\phi(\mathbf{Q}^{(k)})\mathbf{S}^{(k-1)}}
        _{\text{\scriptsize cross-chunk (history)}}
        \;+\; \underbrace{\bigl(\phi(\mathbf{Q}^{(k)})\phi(\mathbf{K}^{(k)})^\top
        \odot \mathbf{M}_C\bigr)\mathbf{V}^{(k)}}
        _{\text{\scriptsize intra-chunk (parallel)}}$\\[10pt]
       $\mathbf{S}^{(k)} = \mathbf{S}^{(k-1)}
        + \phi(\mathbf{K}^{(k)})^\top\mathbf{V}^{(k)}$};
\node[tfnote=tfGreen] (cost) at (\tfrow{18.4}{eq.east})
      {$\mathcal{O}(NC)$ time\\$M = N/C$ sequential steps};

\begin{pgfonlayer}{grouping}
  \node[tfpanel, inner sep=0.7cm, fit=(s0)(sm)(ca)(cao)(cco)(eq)(cost)] (panel) {};
\end{pgfonlayer}
\node[tftitle] at ($(panel.north west)+(0.1,0.15)$)
      {C. Chunk-wise Parallel Form (SSD-style)};

\end{tikzpicture}}
